\chapter{Ingresos presupuestarios}

\marginal{13.4 de cada 100, contra 19.2}
La Ley de Ingresos aprobada para 2025 estima un total de \textbf{9,302,015.8
millones de pesos}, 2.6\% más que 2024 en términos nominales y \textbf{1.6\% menos
en términos reales}. El agregado cae, y sin embargo el gasto financiado con
recursos propios sube: lo que se contrae es el financiamiento.

De cada 100 pesos que el gobierno federal podrá gastar en 2025, \textbf{13.4
provienen de endeudamiento}. En 2024 eran 19.2. Es la caída más grande de esa
proporción en el periodo que cubre esta serie, y es el hecho central del paquete.

\begin{table}[htbp]
\centering
\caption{Ingresos por renglón del artículo 1o., aprobado 2024 contra aprobado 2025 (mdp)}
\label{tab:C2_1}
\small
\begin{tabular}{llSSSSSS}
\toprule
{Clave} & {Concepto} & {2024} & {2025} & {Dif. nom.} & {Dif. real} & {Var. real \%} & {Cámara} \\
\midrule
\num{1} & Impuestos & \num{4942030.3} & \num{5297812.9} & \num{355782.6} & \num{143275.3} & \num{2.8} & \num{0.0} \\
\num{1.11} & Impuestos Sobre los Ingresos & \num{2709899.5} & \num{2859575.1} & \num{149675.6} & \num{33149.9} & \num{1.2} & \num{0.0} \\
\num{1.13} & Impuestos Sobre la Producción, el Consumo y las Transacciones & \num{2037929.8} & \num{2197425.8} & \num{159496.0} & \num{71865.0} & \num{3.4} & \num{0.0} \\
\num{1.14} & Impuestos al Comercio Exterior & \num{101976.3} & \num{151789.7} & \num{49813.4} & \num{45428.4} & \num{42.7} & \num{0.0} \\
\num{1.17} & Accesorios de impuestos & \num{84283.0} & \num{81497.4} & \num{-2785.6} & \num{-6409.8} & \num{-7.3} & \num{0.0} \\
\num{1.18} & Otros impuestos & \num{7811.0} & \num{7140.1} & \num{-670.9} & \num{-1006.8} & \num{-12.4} & \num{0.0} \\
\num{2} & Cuotas y Aportaciones de Seguridad Social & \num{535254.7} & \num{603077.9} & \num{67823.2} & \num{44807.2} & \num{8.0} & \num{0.0} \\
\num{2.22} & Cuotas para la Seguridad Social & \num{535254.7} & \num{603077.9} & \num{67823.2} & \num{44807.2} & \num{8.0} & \num{0.0} \\
\num{3} & Contribuciones de Mejoras & \num{36.5} & \num{38.8} & \num{2.3} & \num{0.7} & \num{1.9} & \num{0.0} \\
\num{4} & Derechos & \num{59091.4} & \num{137500.5} & \num{78409.1} & \num{75868.2} & \num{123.1} & \num{0.0} \\
\num{5} & Productos & \num{8641.6} & \num{13707.1} & \num{5065.5} & \num{4693.9} & \num{52.1} & \num{0.0} \\
\num{6} & Aprovechamientos & \num{193877.0} & \num{223166.3} & \num{29289.3} & \num{20952.6} & \num{10.4} & \num{0.0} \\
\num{7} & Ingresos por Ventas de Bienes, Prestación de Servicios y Otros Ingresos & \num{1312289.4} & \num{1500579.0} & \num{188289.6} & \num{131861.2} & \num{9.6} & \num{0.0} \\
\num{9} & Transferencias, Asignaciones, Subsidios y Subvenciones, y Pensiones y & \num{277774.3} & \num{279766.8} & \num{1992.5} & \num{-9951.8} & \num{-3.4} & \num{0.0} \\
\num{0} & Ingresos Derivados de Financiamientos & \num{1737050.6} & \num{1246366.5} & \num{-490684.1} & \num{-565377.3} & \num{-31.2} & \num{0.0} \\
\bottomrule
\end{tabular}
\par\vspace{2pt}\footnotesize Fuente: elaboración propia del ITED con la Ley de Ingresos aprobada de 2024 y 2025. Las diferencias en millones de pesos se reportan dos veces, nominal y real; la real está en pesos de 2025 con el deflactor de 4.3\% que declara el CGPE.
\end{table}


\section{Qué crece y qué no}

Los impuestos pasan de 4,942,030.3 a 5,297,812.9 millones de pesos, un aumento
real de \textbf{2.8\%}. Las cuotas de seguridad social crecen más, \textbf{8.0\%
real}, de 535,254.7 a 603,077.9. Los ingresos derivados de financiamientos caen
\textbf{31.2\% real}, de 1,737,050.6 a 1,246,366.5.

\marginal{elasticidad implícita de 1.27}
Con un crecimiento real implícito de 2.21\% en el marco macroeconómico y un
aumento real de la recaudación tributaria de 2.8\%, \textbf{la elasticidad
implícita del paquete es de 1.27}. Es un supuesto exigente sin ser extravagante:
supone que la recaudación crece un cuarto más rápido que la economía sin cambios
en la ley. El paquete no lo justifica y no publica el desglose que permitiría
evaluarlo.

\section{El renglón que hay que mirar dos veces}

Los derechos pasan de 59,091.4 a 137,500.5 millones de pesos, \textbf{un aumento
nominal de 132.7\%}. Es la variación más grande de todo el artículo primero y
merece explicación, porque un renglón que se duplica largo en un año rara vez es
recaudación nueva.

El desglose la da en parte. Los derechos por uso y aprovechamiento de bienes de
dominio público suben de 46,414.7 a 78,806.7. Y los derechos por prestación de
servicios pasan de 12,676.7 a 58,693.8, un aumento de 46,017.1 millones que se
concentra en dos dependencias: \textbf{Gobernación pasa de 150.4 a 28,227.4
millones de pesos}, y Hacienda de una cifra menor a 9,656.9.

Un salto de ese tamaño en un renglón de derechos por servicios no es
comportamiento de una base gravable; \textbf{tiene la forma de una
reclasificación o de una tarifa nueva}. El paquete no lo explica: no hay
exposición de motivos del proyecto de presupuesto disponible y la iniciativa de
ingresos no comenta el renglón. Se deja registrado como el cambio de composición
más grande del artículo primero, con su explicación pendiente de fuente.

\section{Sin miscelánea fiscal}

\textbf{El paquete 2025 no incluye una iniciativa de reformas fiscales.} La
palabra no aparece una sola vez en la iniciativa de Ley de Ingresos. La
consolidación de dos puntos del PIB se plantea, por tanto, íntegramente por el
lado del gasto y del financiamiento, sin tocar la estructura tributaria.

Es una decisión con consecuencia de mediano plazo y el propio marco oficial la
muestra: en el escenario a 2030 los ingresos presupuestarios \textbf{bajan} de
22.3 a 21.8\% del PIB, y los tributarios de 14.6 a 14.4. El paquete que consolida
no propone la reforma que haría sostenible la consolidación.

\section{Comparación estructural}

\begin{table}[htbp]
\centering
\caption{Comparaciones estructurales del ITED, series con base consistente}
\label{tab:CR_1}
\footnotesize
\setlength{\tabcolsep}{4pt}
\begin{tabular}{cL{5.2cm}ccccc}
\toprule
{Cap.} & {Razón} & {2021} & {2022} & {2023} & {2024} & {2025} \\
\midrule
\num{2} & {Impuestos / gastos obligatorios con pensiones} & \num{0.671} & \num{0.674} & {---} & \num{0.674} & \num{0.689} \\
\num{2} & {Impuestos + cuotas IMSS / gastos obligatorios con pensiones} & \num{0.743} & \num{0.745} & {---} & \num{0.747} & \num{0.768} \\
\num{4} & {Pensiones y jubilaciones + costo financiero / impuestos} & {---} & {---} & {RETIRADO} & \num{0.559} & \num{0.571} \\
\num{5} & {Pensiones IMSS / cuotas IMSS (serie aprobada)} & {---} & \num{1.545} & \num{1.593} & \num{1.627} & \num{1.603} \\
\num{6} & {Gasto en salud del IMSS / cuotas IMSS} & \num{0.85} & {---} & {---} & \num{0.851} & \num{0.798} \\
\num{6} & {Pensiones IMSS+ISSSTE / salud IMSS+ISSSTE} & \num{2.11} & {---} & {---} & \num{2.282} & \num{2.412} \\
\num{7} & {Función Educación / pensiones y jubilaciones} & {---} & {---} & \num{0.709} & \num{0.689} & \num{0.662} \\
\num{7} & {Función Educación / pensiones totales (con no contributivas)} & {---} & {---} & \num{0.558} & \num{0.518} & \num{0.501} \\
\num{12} & {Costo financiero / impuestos} & {---} & \num{0.20} & \num{0.23} & \num{0.256} & \num{0.262} \\
\num{12} & {SHRFSPF / impuestos, en años de recaudación} & {---} & \num{3.6} & \num{3.4} & \num{3.4} & \num{3.51} \\
\bottomrule
\end{tabular}
\par\vspace{2pt}\footnotesize Fuente: elaboración propia del ITED; numerador y denominador en las columnas correspondientes. Aprobado contra aprobado en toda la serie. Ninguna línea híbrida. El punto de 2023 de la razón del capítulo 4 se retira porque no se reproduce.
\end{table}


La razón que esta casa sigue para el capítulo de ingresos es \emph{impuestos sobre
gastos obligatorios}, con el numerador en el artículo primero de la Ley de
Ingresos y el denominador en el Anexo 3 del decreto. En 2025 vale \textbf{0.690},
contra 0.674 en 2024 y 0.671 en 2021. Con las cuotas de seguridad social en el
numerador, \textbf{0.768} contra 0.747.

La lectura es que la recaudación cubre una porción creciente de los compromisos
inevitables, y que esa mejora viene de las cuotas más que de los impuestos: la
segunda línea sube 2.1 puntos entre 2024 y 2025 mientras la primera sube 1.6. La
razón sigue por debajo de uno, que es lo que importa: \textbf{ni con las cuotas
dentro alcanza la recaudación para pagar lo que ya está comprometido}, y la
diferencia se cubre con deuda.

\clearpage
